\documentclass[11pt]{article}
\usepackage[T1]{fontenc}
\usepackage{lmodern}
\usepackage[margin=1in]{geometry}
\usepackage{amsmath,amssymb,amsthm}
\usepackage[hidelinks]{hyperref}
\newtheorem{theorem}{Theorem}[section]
\newtheorem{proposition}[theorem]{Proposition}
\newtheorem{lemma}[theorem]{Lemma}
\newtheorem{corollary}[theorem]{Corollary}
\theoremstyle{definition}
\newtheorem{definition}[theorem]{Definition}
\theoremstyle{remark}
\newtheorem{remark}[theorem]{Remark}
\title{Evidence-Driven Capacity Change\\Certificates, Reconstruction, and Live Gradient Extensions}
\author{}
\date{4 November 2026}
\begin{document}
\maketitle
\begin{abstract}
A broad structural account of learning would explain how evidence changes
representations across automata and other reconstruction settings. The
concrete question addressed here is narrower. Finite evidence can prove
that a fixed representation is too small without
showing how to enlarge it constructively. This note proves a bounded bridge:
Hankel rank certificates; witness-directed repair and exact reconstruction;
function-preserving gated automaton extensions with explicit loss laws;
noise-robust optimal rank-one births in an explicitly unstructured matrix
relaxation; and reverse-mode compatibility across changing parameter spaces.
Counterexamples separate matrix rank from automaton consistency, reverse
transport from optimizer commutation, and linear realization dimension from
intrinsic nonlinear state complexity. The statements are proved as a synthesis
of established ideas, with no claim of priority or domain-general unification.
\end{abstract}
\tableofcontents
\bigskip
% Standalone-content module. No documentclass or document environment.
% Requires amsmath, amssymb, amsthm, hyperref; theorem environments
% theorem, proposition, lemma, corollary, definition, remark.
% All labels are prefixed ce:. This source is intentionally uncompiled.

\section{Evidence-driven changes of representational capacity}
\label{ce:section}

\paragraph{Scope and status.}
The motivating objective is a structural account of how evidence changes
what counts as a state or representation, potentially connecting automata,
formal-concept analysis, and topological reconstruction. The results here
address a specific obstruction to that broader objective: changing capacity
inside a differentiable learner. They establish no general equivalence of
those reconstruction domains.

The useful question is not simply whether a learner can differentiate its
parameters. It is whether observations can certify the inadequacy of its
current representation, and whether one can then enlarge the representation
without discarding its learned behavior or creating parameters that gradients
cannot activate. Three mathematically different operations must be separated:
\begin{enumerate}
\item certify that a given capacity is insufficient;
\item optimize within a fixed-capacity family;
\item introduce new, behaviorally effective directions and transport the old
      reverse-mode calculation into the enlarged family.
\end{enumerate}
The results below are mathematical derivations and a synthesis of established
ideas, not claims of priority. Hankel realization and categorical automata
learning already explain exact state discovery; low-rank optimization already
studies rank-increasing steps; function-preserving network transformations and
parametric lenses already study their respective interfaces. The contribution
of this section is a precise separation and integration of these ingredients,
including limitations that prevent a false unification.

\subsection{Finite evidence can certify an inadequate state dimension}

Fix a finite alphabet $\Sigma$. A real weighted finite automaton (WFA) with
$n$ states is a tuple
\[
 \mathcal A=(\alpha,(A_a)_{a\in\Sigma},\beta),\qquad
 \alpha\in\mathbb R^{1\times n},\quad A_a\in\mathbb R^{n\times n},
 \quad\beta\in\mathbb R^{n\times1}.
\]
Write $A_\varepsilon=I$ and $A_w=A_{a_1}\cdots A_{a_k}$ for
$w=a_1\cdots a_k$. Its behavior is $f_{\mathcal A}(w)=\alpha A_w\beta$.
For finite prefix and suffix sets $S,E\subseteq\Sigma^*$, its observed
Hankel block is
\[
 H_f[S,E]_{s,e}=f(se),\qquad H_f[S,E]\in\mathbb R^{|S|\times|E|}.
\]
It factors as $H_f[S,E]=R_SO_E$, where the row indexed by $s$ of $R_S$
is $\alpha A_s$ and the column indexed by $e$ of $O_E$ is $A_e\beta$.
Consequently $\operatorname{rank}H_f[S,E]\leq n$. This is a state-dimension
bound, not a statement about the rank of an arbitrary neural-network weight
matrix.

\begin{theorem}[Exact and noise-robust capacity certificates]
\label{ce:certificate}
Let $H_*\in\mathbb R^{m\times n}$ be a true Hankel block and let
$Y\in\mathbb R^{m\times n}$ satisfy
$\|Y-H_*\|_2\leq\epsilon$, where $\|\cdot\|_2$ is the operator norm and
$\epsilon\geq0$ is a justified error bound. Put $p=\min(m,n)$ and take
$0\leq r<p$. Then:
\begin{enumerate}
\item If $\sigma_{r+1}(Y)>\epsilon$, no WFA with at most $r$ states
      realizes the true behavior.
\item Every rank-at-most-$r$ prediction matrix $X$ obeys
\[
 \|H_*-X\|_2\geq(\sigma_{r+1}(Y)-\epsilon)_+,
 \qquad
 \|H_*-X\|_F^2\geq
 \sum_{j=r+1}^{p}(\sigma_j(Y)-\epsilon)_+^2.
\]
Here $a_+=\max(a,0)$.
\item Thus $\sigma_{r+1}(Y)>\epsilon+\tau$ certifies that an $r$-state
      WFA cannot attain operator-norm error at most $\tau$ on this block.
\end{enumerate}
\end{theorem}
\begin{proof}
Singular-value perturbation gives
$\sigma_j(H_*)\geq\sigma_j(Y)-\epsilon$. The best rank-$r$
approximation error in operator norm is $\sigma_{r+1}(H_*)$, and its
squared Frobenius error is $\sum_{j>r}\sigma_j(H_*)^2$. Applying these
bounds to $X$ proves the inequalities. A WFA with at most $r$ states has
block rank at most $r$, which proves the remaining claims.
\end{proof}

The certificate is only as sound as its noise bound. An empirical singular
value is not itself evidence of an additional true state when the error scale
is unknown. Under a specified sampling model, concentration results for
empirical Hankel matrices can supply such bounds
\cite{ce:denis2016}. The theorem uses no sampling assumptions beyond the
stated event $\|Y-H_*\|_2\leq\epsilon$.
The spectral criterion is sufficient for the structured WFA problem. Its
failure is not a proof that a consistent WFA of order $r$ exists: arbitrary
low-rank matrices need not satisfy the repeated-word equalities or transition
conditions of a Hankel realization.

\begin{proposition}[What exhaustive exact evidence eventually reveals]
\label{ce:eventual}
Suppose a nonzero behavior has a minimal $d$-state real WFA. Then
\[
 \operatorname{rank}H_f[\Sigma^{\leq d-1},\Sigma^{\leq d-1}]=d.
\]
Ranks of nested exact observation blocks are nondecreasing and bounded by
$d$. However, an observed plateau below $d$ is not a certificate of
completion when $d$ is unknown.
\end{proposition}
\begin{proof}
Let $V_k$ be the span of reachable rows $\alpha A_w$ with $|w|\leq k$.
If $V_k=V_{k+1}$, then $V_kA_a\subseteq V_k$ for every $a$, so the
reachable span never grows thereafter. In a minimal realization it is all
of $\mathbb R^{1\times d}$. Since $V_0$ is one-dimensional, either full
reachability is attained earlier or its dimension strictly increases at each
step through $d-1$. The analogous argument for columns $A_w\beta$ proves
full observability by length $d-1$. The resulting $R_S$ and $O_E$ have
column and row rank $d$, respectively, so their product has rank $d$.
Nested blocks are submatrices of larger blocks, giving monotonicity.
Without a bound on $d$, an unseen longer word can still expose an additional
direction; monotonicity alone supplies no stopping rule.
\end{proof}


\subsection{From a witnessed table defect to a repaired realization}
\label{ce:exact-interface}

The next result makes the categorical image construction operational. It is
an exact linear instance of the closedness/consistency method underlying
CALF \cite{ce:calf2017}, not a replacement for that literature.
Assume $S$ is prefix-closed, $E$ is suffix-closed, and both contain
$\varepsilon$. All queried words have coherent, exact real values: the same
word appearing in different entries has the same value. Define
\[
 T:\mathbb R^S\to\mathbb R^E,\quad
 Te_s=(f(se))_{e\in E},\qquad
 K_a e_s=(f(sae))_{e\in E},\qquad I=\operatorname{im}T.
\]
Thus the matrix of $T$ is the transpose of the earlier row-indexed Hankel
block. The observed word set is the set of distinct words
$D=SE\cup S\Sigma E$. Missing entries are not assigned zero.

\begin{theorem}[Exact reconstruction and witness-directed rank progress]
\label{ce:repair}
Put $k=\operatorname{rank}T$.
\begin{enumerate}
\item If for every $a\in\Sigma$,
\[
 \operatorname{im}K_a\subseteq I,
 \qquad \ker T\subseteq\ker K_a,
\]
then there is a unique linear operator $D_a:I\to I$ with $D_aT=K_a$.
Starting at $q_0=Te_\varepsilon$ and reading output coordinate
$\varepsilon$ yields a $k$-dimensional WFA that reproduces every queried
word in $D$. Its minimal realization dimension is exactly $k$.
\item If $\operatorname{im}K_a\not\subseteq I$, some $s\in S$ has
$K_ae_s\notin I$. Appending the prefix $sa$ increases the rank of the
unshifted table by exactly one.
\item If $\ker T\not\subseteq\ker K_a$, choose $v\in\ker T$ and
$e\in E$ with $(K_av)_e\ne0$. Appending the suffix $ae$ increases
that rank by exactly one.
\end{enumerate}
Both repairs preserve the respective word closures. Their new unshifted
entries have already been queried in $S\Sigma E$; completing the enlarged
shifted blocks can require additional queries.
\end{theorem}
\begin{proof}
The kernel condition makes the assignment $D_a(Tv)=K_av$ well-defined;
the image condition makes its values lie in $I$. Surjectivity onto $I$
gives uniqueness. For a word $w=a_1\cdots a_j$, write
$D_w=D_{a_j}\cdots D_{a_1}$, so states evolve in reading order.
Prefix closure and induction give $D_sq_0=Te_s$ for every $s\in S$.
Let $o:I\to\mathbb R$ be coordinate $\varepsilon$. Suffix closure gives
\[
 oD_e(q)=q_e\quad(e\in E,\ q\in I).
\]
Indeed the empty suffix is immediate; for $e=av$ with $v\in E$, induction
reduces the assertion on $q=Te_s$ to
$(D_aTe_s)_v=f(sav)=(Te_s)_{av}$, and these $q$ span $I$.
Consequently the realized values on $se$ and $sae$ are respectively
$(Te_s)_e$ and $(K_ae_s)_e$. Choosing a basis on $I$, with row transitions
$A_a=D_a^T$, puts this realization in the earlier WFA convention. Its
observed block has rank $k$, so no smaller realization has its behavior.
In the zero-rank case, the reconstructed hypothesis is the zero series,
with a zero-dimensional realization
or, if required by a convention, a one-state zero realization.

For the second claim, if every column of $K_a$ lay in $I$, its image
would lie there. The chosen independent column consists of the values
$f(sae')$ and becomes the column for $sa$ in the new $T$. The prefix
$sa$ is genuinely new; its proper prefixes already belong to $S$.
For the third claim, the $e$-row of $K_a$ does not annihilate $v$, while
every row of $T$ does. It is therefore independent of the existing rows,
and it is precisely the new row indexed by $ae$. Its proper suffixes
already lie in $E$. Adding one independent column or row raises rank by
exactly one.
\end{proof}

\begin{corollary}[Finite exact repair and the meaning of a counterexample]
\label{ce:repair-termination}
If the true target has some WFA of dimension at most $N$, repeated
witness-directed repairs, with each enlarged query set completed, make at
most $N-k$ repairs before all operators descend. The resulting hypothesis
fits its finite query set, not necessarily the global target.
If a tested word $w$ refutes an already repaired rank-$k$ hypothesis,
append all prefixes of $w$ to $S$ and complete the new queries. Either
the table rank already exceeds $k$, or it has a defect whose repair raises
the rank beyond $k$.
\end{corollary}
\begin{proof}
Every true observation table has rank at most $N$, and every repair
strictly increases rank. For the counterexample claim, suppose the enlarged
table instead has the same rank and no defect. Since $E$ is unchanged,
the old image is a subspace of the new image of the same dimension, hence
they are equal. The two descended transition operators agree on the old
spanning columns because their shifted entries agree. Their start and
output maps also agree. Thus the new and old hypotheses have the same
behavior. Reconstruction at the new table must, however, reproduce the
observed value at $w$, contradicting the counterexample.
\end{proof}

The query semantics are explicit: for an enlarged context $(S',E')$,
query only the previously unknown distinct words in
$S'E'\cup S'\Sigma E'$. No equivalence oracle is needed to perform a
witness repair. Finding a counterexample, or knowing that none exists, is
a separate experimental obligation. Noisy kernel or image residuals can
suggest queries but are not exact rank-progress proofs. The noisy
certificate of Theorem~\ref{ce:certificate} must be checked on the
resulting block using a justified error bound for that block.

\begin{theorem}[A repaired-hypothesis path with empirical no-forgetting]
\label{ce:reconstructed-gate}
Let $C\subseteq D$ be finite coherent exact observation sets. Suppose an
$n$-state model $f$ fits $C$, and a repaired table constructs an $m$-state
model $h$ fitting $D$ by Theorem~\ref{ce:repair}. There is a smooth
one-parameter family of $(n+m)$-state WFAs with behavior
\[
 f_t=(1-t)f+th.
\]
It has the following properties:
\begin{enumerate}
\item $f_0=f$ on all words.
\item Every old observed value on $C$ is preserved for all $t$.
\item For positive weights on $D$ and squared loss,
$\mathcal L_D(f_t)=(1-t)^2\mathcal L_D(f)$.
If $f$ fails the refined evidence, the gate derivative at zero is
$-2\mathcal L_D(f)<0$, and $t=1$ reaches zero refined loss.
\item The endpoint at $t=1$ can be compressed to the $m$-state model
$h$ without changing any word output.
\end{enumerate}
If the refined table has rank $m>n$, its evidence proves that the old
capacity is insufficient. The temporary dimension bound $n+m$ is an
upper bound, not a minimal extension theorem.
\end{theorem}
\begin{proof}
For $f=(\alpha,(A_a),\beta)$ and
$h=(\gamma,(B_a),\delta)$, take
\[
 \alpha_t=((1-t)\alpha,t\gamma),\quad
 M_{a,t}=\operatorname{diag}(A_a,B_a),\quad
 \beta_t=\binom\beta\delta.
\]
Block multiplication gives the wordwise affine combination. Both endpoint
models agree with the same labels on $C$, proving retention there.
Since $h$ fits $D$, every refined residual is $(1-t)$ times the old
residual. At $t=1$ the old block has zero initial weight and can be
removed. The rank obstruction follows from Theorem~\ref{ce:certificate}
with zero noise.
\end{proof}

This theorem gives an explicit safe one-dimensional path with endpoints
held fixed. It is not a no-forgetting guarantee for unconstrained updates
of all enlarged parameters. Unlike finite-support memorization, its new
endpoint comes from a repaired realization whose dimension is the evidence
image dimension. Nevertheless its agreement on unqueried words remains
unverified, and noisy exact interpolation remains unjustified.

\subsection{A robust, optimal rank-one birth in the matrix relaxation}

This subsection concerns matrices, not automatically WFA realizations. Its
strength is that a new direction is selected from evidence and yields a
quantified improvement even for the unknown noise-free table.

Let $X\in\mathbb R^{m\times n}$ have rank $r<\min(m,n)$, and write a
compact SVD as $X=U\Lambda V^T$. Put $P=UU^T$, $Q=VV^T$. The normal
space to the rank-$r$ manifold at $X$ is
\[
 \mathcal N_X=\{Z:PZ=0,\ ZQ=0\}.
\]
Given an observed matrix $Y$, define its normal innovation
\[
 N=(I-P)(Y-X)(I-Q)=(I-P)Y(I-Q),\qquad\rho=\|N\|_2.
\]
When $\rho>0$, choose unit vectors $u\perp\operatorname{im}U$ and
$v\perp\operatorname{im}V$ with $u^TNv=\rho$.

\begin{theorem}[Noise-robust, minimax normal birth]
\label{ce:robust-birth}
Suppose $\|Y-H_*\|_2\leq\epsilon$ and $\rho>\epsilon$. Define
$t=\rho-\epsilon$ and
\[
 X_+=X+tuv^T,\qquad
 \mathcal L_T(Z)=\tfrac12\|Z-T\|_F^2.
\]
Then:
\begin{enumerate}
\item $\operatorname{rank}X_+=r+1$.
\item The unknown true table loss improves by the certified amount
\[
 \mathcal L_{H_*}(X)-\mathcal L_{H_*}(X_+)
 \ \geq\ \tfrac12(\rho-\epsilon)^2.
\]
\item This guaranteed improvement is optimal among all nonnegative step
      sizes and unit normal rank-one atoms under the unstructured
      operator-norm uncertainty ball. Precisely,
\[
 \sup_{\substack{a,b\ {\mathrm{unit}}\\ Pa=0,\ Qb=0}}\ \sup_{s\geq0}
 \inf_{\|T-Y\|_2\leq\epsilon}
 \bigl[\mathcal L_T(X)-\mathcal L_T(X+sab^T)\bigr]
 =\tfrac12(\rho-\epsilon)_+^2.
\]
The displayed construction attains the supremum when $\rho>\epsilon$.
\end{enumerate}
\end{theorem}
\begin{proof}
In orthonormal bases extending $U,u$ and $V,v$, the nonzero part of
$X+tuv^T$ is block diagonal with blocks $\Lambda$ and $t$. This proves
its rank.
For any unit $a,b$ and any $s\geq0$, direct expansion gives
\[
 \mathcal L_T(X)-\mathcal L_T(X+sab^T)
 =s\langle T-X,ab^T\rangle_F-\tfrac12s^2.
\]
The least possible value of $\langle T-Y,ab^T\rangle_F$ over
$\|T-Y\|_2\leq\epsilon$ is $-\epsilon$: the lower bound follows from
operator/nuclear norm duality and is attained by
$T-Y=-\epsilon ab^T$. For a normal atom,
$\langle Y-X,ab^T\rangle_F=a^TNb$. Consequently its worst-case
improvement is exactly
\[
 s(a^TNb-\epsilon)-\tfrac12s^2.
\]
The greatest $a^TNb$ is $\rho$. Maximization of the remaining scalar
quadratic over $s\geq0$ gives $s=(\rho-\epsilon)_+$ and value
$\tfrac12(\rho-\epsilon)_+^2$. The actual $H_*$ is in the uncertainty
ball, which proves the claimed guarantee.
\end{proof}

The minimax qualification matters: the uncertainty set in the equality above
contains all matrices in the norm ball, including matrices that are not
Hankel. The same bound remains valid if the true table is Hankel, but
optimality relative to a smaller structured uncertainty set is not claimed.
Nor is this an out-of-sample generalization bound.

\begin{corollary}[A capacity certificate forces birth after fixed-rank stationarity]
\label{ce:stationary-birth}
Assume $X$ is stationary for
$\frac12\|X-Y\|_F^2$ on the smooth manifold of rank-exactly-$r$ matrices.
Then $Y-X=N$ and
\[
 \rho=\|Y-X\|_2\geq\sigma_{r+1}(Y).
\]
In particular, $\sigma_{r+1}(Y)>\epsilon$ forces the birth in
Theorem~\ref{ce:robust-birth} and guarantees true loss decrease at least
$\frac12(\sigma_{r+1}(Y)-\epsilon)^2$.
\end{corollary}
\begin{proof}
The tangent space consists of matrices of the form $UK^T+LV^T$.
Orthogonality of $X-Y$ to every such tangent vector implies
$U^T(X-Y)=0$ and $(X-Y)V=0$. Hence $Y-X$ is its own normal projection.
Every rank-$r$ approximation has operator-norm error at least
$\sigma_{r+1}(Y)$.
\end{proof}

At an arbitrary iterate a rank certificate does not imply $N\ne0$:
some error may first require movement tangent to the existing rank stratum.
For example $X=\left(\begin{smallmatrix}1&0\\0&0\end{smallmatrix}\right)$
and $Y=\left(\begin{smallmatrix}1&1\\1&0\end{smallmatrix}\right)$ have
$\operatorname{rank}Y=2$ but $N=0$ at $X$. This is why fixed-capacity
optimization and certified expansion must not be conflated.


Normal descent and intrinsic rank insufficiency remain different claims.
Even stationarity together with $\rho>\epsilon$ does not by itself certify
true rank greater than $r$ under noise. For example, let
\[
 r=1,\quad X=\operatorname{diag}(0.05,0),\quad
 Y=\operatorname{diag}(0.05,1),\quad
 H_*=\operatorname{diag}(0,1),\quad\epsilon=0.1.
\]
The point $X$ is stationary on the rank-one manifold, and $\rho=1$,
but the true table has rank one. The birth improves a poor stationary
representation; it does not prove that all rank-one representations fail.
The separate criterion $\sigma_{r+1}(Y)>\epsilon$ is sufficient.
At a best rank-$r$ truncated SVD, $\rho=\sigma_{r+1}(Y)$, but that
equality need not hold at other stationary points.

\begin{proposition}[Dead symmetric padding and live asymmetric padding]
\label{ce:factor-padding}
Let $X=AB^T$, where $A\in\mathbb R^{m\times r}$ and
$B\in\mathbb R^{n\times r}$ have full column rank, and let $F$ be a
continuously differentiable function of $X$ with gradient $G$.
For the enlarged factorization
$\Phi(A_+,B_+)=A_+B_+^T$:
\begin{enumerate}
\item At $A_+=[A,0]$, $B_+=[B,0]$, both new-column gradients are zero.
      Gradient descent initialized there cannot activate either new column.
\item At $A_+=[A,u]$, $B_+=[B,0]$, the represented matrix is still $X$,
      but the new-column gradient in $B_+$ is $G^Tu$.
\item For $F(Z)=\frac12\|Z-Y\|_F^2$ and the normal singular vectors
      above, the path $B_+(s)=[B,sv]$ has
\[
 \left.\frac{d}{ds}F(A_+B_+(s)^T)\right|_{s=0}=-\rho.
\]
      Its matrix is $X+suv^T$ and has rank $r+1$ for every $s\ne0$.
\end{enumerate}
\end{proposition}
\begin{proof}
The factor gradients are $\nabla_{A_+}(F\circ\Phi)=GB_+$ and
$\nabla_{B_+}(F\circ\Phi)=G^TA_+$. These give the first two claims,
and the all-zero new columns remain zero under each gradient-descent step.
For the last claim, $u^TGv=-u^T(Y-X)v=-\rho$; orthogonality to the old
column and row spaces proves the rank statement.
\end{proof}

\paragraph{Worked numerical example.}
Take $Y=\operatorname{diag}(2,0.8)$, $X=\operatorname{diag}(2,0)$,
$r=1$, and $\epsilon=0.1$. The old rank-one point is stationary,
$u=v=(0,1)^T$, and $\rho=0.8$. The robust step gives
$X_+=\operatorname{diag}(2,0.7)$ and guarantees a reduction of at least
$0.245$ in the unknown true half-squared Frobenius loss. This is attained
by the admissible target $H_*=\operatorname{diag}(2,0.7)$.
The observed loss drops from $0.32$ to $0.005$.
Appending two zero factor columns would instead produce zero gradient in
both new columns. Appending $u$ to one factor and a zero column to the other
preserves the starting matrix while making the scalar gate derivative
$-0.8$.

\paragraph{An explicit obstruction to interpreting the step as an automaton.}
For a unary alphabet and $S=E=\{\varepsilon,a,a^2\}$, consider
\[
 X=\begin{pmatrix}1&0&0\\0&0&0\\0&0&0\end{pmatrix},\qquad
 Y=\begin{pmatrix}1&0&1\\0&1&0\\1&0&0\end{pmatrix}.
\]
Both are valid finite Hankel blocks. At $X$, the top normal innovation is
$uv^T=e_2e_2^T$. But $X+te_2e_2^T$ is not Hankel for $t\ne0$:
its $(a,a)$ entry differs from its $(\varepsilon,a^2)$ and
$(a^2,\varepsilon)$ entries, although all three must evaluate $f(a^2)$.
Projecting this atom onto the Hankel subspace gives
\[
 \tfrac13\begin{pmatrix}0&0&1\\0&1&0\\1&0&0\end{pmatrix},
\]
which has rank three, not one. Thus imposing the linear observation
constraints can destroy the very rank-one structure used in the birth
argument. A genuine low-order automaton extension must also respect
repeated-word identities and transition consistency. Constrained Hankel
completion is therefore relevant prior art, rather than an optional
implementation detail \cite{ce:balle2012}.

\subsection{Reverse-mode compatibility is not optimizer commutation}

\begin{proposition}[The correct compatibility equation]
\label{ce:pullback}
Let $P\subseteq\mathbb R^d$ and $P'\subseteq\mathbb R^{d'}$ be open,
let $e:P\to P'$ be smooth, and let $F:P\times\mathcal X\to\mathbb R^k$
and $F':P'\times\mathcal X\to\mathbb R^k$ be smooth in their parameter
arguments. Suppose
\[
 F(p,x)=F'(e(p),x)\quad\hbox{for every }p,x.
\]
Then parameter reverse derivatives satisfy, for every output covector
$\bar y\in\mathbb R^k$,
\[
 D_pF(p,x)^T\bar y
 =De(p)^T D_{p'}F'(e(p),x)^T\bar y.
\]
For any common differentiable loss, the corresponding parameter-loss
covectors therefore obey $dL_p=De(p)^T dL'_{e(p)}$.
\end{proposition}
\begin{proof}
Differentiate $F(-,x)=F'(-,x)\circ e$ and transpose the chain rule.
Apply the same argument after composing with the loss.
\end{proof}

This is the relevant reverse-derivative/lens compatibility. It transports
old covectors correctly without requiring the enlarged optimizer to stay in
the old parameter subspace. Turning covectors into update vectors requires
additional choices of metric or optimizer, as in parametric-lens treatments
\cite{ce:lenses2024}. No Euclidean-gradient invariance under arbitrary
non-isometric coordinate changes is asserted.

\begin{theorem}[Exact failure of optimizer commutation enables escape]
\label{ce:noncommutation}
Suppose specifically that $e(p)=Jp+c$, where
$J\in\mathbb R^{d'\times d}$ has $J^TJ=I$. Let $L=L'\circ e$.
At a point $p$, write
\[
 g=\nabla L(p),\qquad G=\nabla L'(e(p)),\qquad
 n=(I-JJ^T)G.
\]
Then $G=Jg+n$, $J^Tn=0$, and for equal-step Euclidean gradient descent,
\[
 \bigl(e(p)-\eta G\bigr)-e(p-\eta g)=-\eta n.
\]
Thus the two updates commute at this point, for $\eta\ne0$, if and only
if $n=0$. If $n\ne0$, the new normal direction is a strict first-order
loss-descent direction:
\[
 \left.\frac{d}{dt}L'(e(p)-tn)\right|_{t=0}=-\|n\|^2.
\]
If $\nabla L'$ is $K$-Lipschitz on the relevant line segment, $K>0$,
then moving only in that direction with $0<t<2/K$ gives
\[
 L'(e(p)-tn)\leq L'(e(p))
 -t(1-Kt/2)\|n\|^2.
\]
\end{theorem}
\begin{proof}
Proposition~\ref{ce:pullback} gives $g=J^TG$. Orthogonal decomposition
then gives the formula for $G$, and the affine form of $e$ gives the exact
update identity. Finally $\langle G,n\rangle=\|n\|^2$; the descent
lemma proves the finite-step bound.
\end{proof}

Insisting on optimizer commutation would therefore remove precisely the
normal component by which the larger learner can depart from the smaller
one. Conversely, parameter-space departure alone does not certify increased
functional capacity: redundant parametrizations can contain normal
parameters without representing genuinely new behaviors. A separate
behavioral witness such as Theorem~\ref{ce:certificate} is needed.

\subsection{A genuine automaton extension with a live gradient}

The matrix construction is economical but may be dynamically inconsistent.
The following construction is dynamically consistent but can be expensive.
Together they separate the real problem from a merely formal bridge.

\begin{proposition}[The WFA version of the zero-padding obstruction]
\label{ce:wfa-dead}
Pad an $r$-state WFA by $q$ extra coordinates as
\[
 \alpha'=(\alpha,0),\quad
 A'_a=\begin{pmatrix}A_a&0\\0&C_a\end{pmatrix},\quad
 \beta'=\binom\beta0,
\]
where the $C_a$ are arbitrary. For every word, the derivative of its
prediction with respect to every new initial/output coordinate, every
cross-block transition coordinate, and every coordinate of the $C_a$ is
zero. Consequently all these coordinates have zero gradient for every
differentiable loss depending on finitely many word predictions. Pure
gradient descent starting on this padded submanifold cannot activate the
new states.
\end{proposition}
\begin{proof}
The prediction derivative is a sum of an initial-vector variation, an
output-vector variation, and terms
$\alpha'A'_u(\delta A'_a)A'_v\beta'$.
A variation in a new initial coordinate must terminate through the zero new
output vector; a new output variation must begin through the zero new
initial vector. A single cross-block transition cannot both enter and leave
the new block. A new-block internal variation cannot be reached from the
old block. All the stated derivatives vanish. The same structure persists
after an update of the old coordinates alone.
\end{proof}


\begin{theorem}[An economical one-state birth, conditional on correlation]
\label{ce:one-state}
Let $f=(\alpha,(A_a),\beta)$ be a minimal $r$-state WFA, $r\geq1$,
over a nonempty alphabet. Choose a letter $a_0$ and a real
$\lambda\notin\operatorname{spec}(A_{a_0})$. Define
\[
 \alpha^+=(\alpha,1),\quad
 A_a^+=\operatorname{diag}(A_a,\lambda),\quad
 \beta^+(t)=\binom\beta t.
\]
Then $f_t(w)=f(w)+t\lambda^{|w|}$, and every $t\ne0$ gives a behavior
of minimal dimension exactly $r+1$. For finitely many weighted observations,
put $e_i=y_i-f(w_i)$,
$C_\lambda=\sum_i\omega_i e_i\lambda^{|w_i|}$ and
$G_\lambda=\sum_i\omega_i\lambda^{2|w_i|}$. If $C_\lambda\ne0$,
the gate derivative is $-C_\lambda$, and its exact squared-loss minimizer
is $t=C_\lambda/G_\lambda$, with loss reduction
$C_\lambda^2/(2G_\lambda)>0$.
\end{theorem}
\begin{proof}
Block multiplication proves the behavior identity. Let
$p(z)=\det(zI-A_{a_0})$. Cayley--Hamilton gives
\[
 \alpha^+p(A_{a_0}^+)=(0,p(\lambda)),\qquad
 p(A_{a_0}^+)\beta^+(t)=\binom0{t p(\lambda)}.
\]
These are linear combinations of reachable rows and observable columns,
respectively. Since $p(\lambda)\ne0$, they isolate the new coordinate
when $t\ne0$. Subtracting that coordinate from every reachable row
$(\alpha A_w,\lambda^{|w|})$ and every observable column
$\binom{A_w\beta}{t\lambda^{|w|}}$ recovers the old reachable and observable
spans. Minimality of the original model makes both enlarged spans have
dimension $r+1$. Finite bases of those spans yield a Hankel block of rank
$r+1$, proving minimality. Finally the gate loss is
$L(t)=L(0)-C_\lambda t+G_\lambda t^2/2$.
\end{proof}

This one-state construction need not preserve old observed values after
the gate moves. With one fresh nonzero error and all prior samples exactly
fitted, any positive admissible $\lambda$ gives $C_\lambda\ne0$.
For a batch, cancellation can make every such gate ineffective: the
polynomial $C_\lambda$ vanishes identically when the weighted residual
sums vanish within each word-length class. Even allowing one scalar
$\lambda_a$ per letter cannot distinguish $ab$ from $ba$; opposite
residuals on those words cancel for all scalar-state branches. This is a
limitation of the stated additive branch, not a proof that all coupled
$(r+1)$-state architectures fail.

A small coupled example shows what a more structured repair can accomplish.
For the unary old behavior $f(a^j)=1$, keep the old observations
$f(\varepsilon)=f(a)=1$ and add the label $y(a^2)=2$. Its observed
Hankel block is $\left(\begin{smallmatrix}1&1\\1&2\end{smallmatrix}\right)$,
which forces dimension two. The two-state family
\[
 \alpha=(1,0),\quad\beta=(1,0)^T,\quad
 A(t)=\begin{pmatrix}1&1\\t&0\end{pmatrix}
\]
preserves the whole old constant behavior at $t=0$, preserves the two old
observations for every $t$, and has $f_t(a^2)=1+t$. Its fresh-observation
loss is $(t-1)^2/2$, with nonzero gate derivative at zero. The same
$2\times2$ Hankel block has determinant $t$, so $t\ne0$ activates
exactly one genuinely necessary state. This example is constructive, not a
universal one-state no-forgetting theorem.

\begin{lemma}[A finite behavioral residual has an explicit realization]
\label{ce:trie}
Let $D\subseteq\Sigma^*$ be finite and nonempty, with one real coefficient
$c_w$ for each distinct word $w\in D$. Let
$g(w)=c_w$ for $w\in D$ and $g(w)=0$ otherwise.
Then $g$ has a WFA with $q=|\operatorname{Pref}(D)|$ states, including
the empty prefix. Its initial row is the basis row $e_\varepsilon^T$;
its transition matrix $B_a$ has an entry $1$ from $p$ to $pa$ when both
are prefixes of $D$, and zero entries elsewhere; its output entry at $p$
is $c_p$ if $p\in D$ and $0$ otherwise.
\end{lemma}
\begin{proof}
After reading $w$, the initial row is exactly the basis row at $w$ if $w$
is a prefix of a word in $D$, and is zero otherwise. Multiplication by the
output column returns exactly the specified coefficient, including when
one observed word is a prefix of another.
\end{proof}

\begin{theorem}[Behavior-preserving, gradient-active automaton expansion]
\label{ce:residual-extension}
Let $f_\theta$ be an $r$-state WFA. Fix a finite nonempty labeled set
$D$ with labels $y_w$ and weights $\omega_w>0$, and fix the current
parameter $\theta_0$. Define
\[
 \mathcal L(\theta)=\tfrac12\sum_{w\in D}\omega_w
       (f_\theta(w)-y_w)^2,
 \qquad
 g(w)=\begin{cases}y_w-f_{\theta_0}(w),&w\in D,\\0,&w\notin D.
 \end{cases}
\]
Hold this residual function fixed during the following expansion. Let
$(\gamma,(B_a),\delta)$ be its trie realization, or any exact smaller
realization. Define
\[
 \widetilde\alpha=(\alpha,t\gamma),\qquad
 \widetilde A_a=\operatorname{diag}(A_a,B_a),\qquad
 \widetilde\beta=\binom\beta\delta.
\]
Then:
\begin{enumerate}
\item This is an $(r+q)$-state WFA with behavior
      $\widetilde f_{\theta,t}=f_\theta+tg$.
\item At $t=0$ it preserves $f_\theta$ on every word, not just observed
      words. For every $t$, it preserves the current behavior outside $D$
      if $\theta=\theta_0$ is held fixed.
\item With $\theta=\theta_0$ held fixed,
\[
 \widetilde{\mathcal L}(\theta_0,t)
 =(1-t)^2\mathcal L(\theta_0),\qquad
 \partial_t\widetilde{\mathcal L}(\theta_0,0)
 =-2\mathcal L(\theta_0).
\]
      Thus any nonzero residual yields a live gate gradient; every
      $0<t<2$ strictly improves the empirical loss, and $t=1$
      interpolates the labels.
\item At the embedding $e(\theta)=(\theta,0)$, the old parameter-loss
      gradient is exactly preserved. At the reference point $\theta_0$,
      the additional gate component is nonzero whenever
      $\mathcal L(\theta_0)>0$. No such claim is made for arbitrary
      $\theta$ against the frozen residual.
\item Suppose $SE\subseteq D$ for some finite $S,E$, and the label block
      $Y_{s,e}=y_{se}$ has rank greater than $r$. Then there is $\eta>0$
      such that $\widetilde f_{\theta_0,t}$ has minimal state dimension
      greater than $r$ for every $0<t<\eta$.
\end{enumerate}
\end{theorem}
\begin{proof}
Block-diagonal multiplication gives $\widetilde f=f+tg$, proving the
first two claims. On every labeled word the new residual is $(1-t)$ times
the old residual, which proves the loss identity and its derivative.
Since $\widetilde f_{\theta,0}=f_\theta$, the chain rule proves old-gradient
preservation. In the coordinates $(\theta,t)$ this is an isometric affine
embedding, so the nonzero gate component is also an explicit instance of
Theorem~\ref{ce:noncommutation}.

For the last claim, let $X=H_{f_{\theta_0}}[S,E]$, so
$\operatorname{rank}X\leq r$. The new block is
$X(t)=(1-t)X+tY$. Since $\operatorname{rank}Y>r$, some
$(r+1)\times(r+1)$ minor has determinant $p(t)$ with $p(1)\ne0$.
This polynomial is not identically zero and has only finitely many roots.
There is therefore $\eta>0$ such that $p(t)\ne0$ for all
$0<t<\eta$. The new block, hence every realization of its behavior, has
rank greater than $r$ on that interval.
\end{proof}


\begin{corollary}[One fresh word can be repaired without forgetting]
\label{ce:single-word}
For a fresh word $w_*$ of length $j$, absent from the old sample set,
there is an extension with at most $j+1$ additional states and behavior
$f_t=f+t\mathbf1_{\{w_*\}}$. It preserves every other word output for
all $t$ and fits the fresh label at $t=y_*-f(w_*)$.
\end{corollary}
\begin{proof}
Apply the trie construction to the singleton support $\{w_*\}$, which
has $j+1$ prefixes. Its output coefficient is the scalar gate $t$.
\end{proof}

\paragraph{A general differentiable loss.}
The construction is not restricted to squared error. Let $z_0$ be the
vector of current predictions on distinct words in $D$, and let
$\ell:\mathbb R^{|D|}\to\mathbb R$ be differentiable. If
$h=\nabla\ell(z_0)\ne0$, realize the finitely supported residual
$g(w)=-h_w$ on $D$. Then
\[
 \left.\frac{d}{dt}\ell(z_0+t\,g|_D)\right|_{t=0}=-\|h\|^2<0.
\]
A nonzero prediction-space gradient is the necessary assumption here:
positive loss alone does not imply that property for an arbitrary loss.

\paragraph{Why this is an upper bound, not the research problem solved.}
The trie proves that a consistent and live expansion always exists on a
finite dataset. It may merely memorize, use many extra states, and provide
no improvement outside the evidence. Exact realization minimization can
replace the trie by an automaton of dimension $\operatorname{rank}H_g$,
but does not make that residual dimension small. Under noise, interpolation
can overfit even when a separate robust certificate correctly says that
$r$ states are insufficient. One must not describe this theorem as a
minimal or statistically optimal capacity-selection algorithm.

A particularly sharp example is unary $f(a^k)=1$ for all $k$, with
$D=\{\varepsilon,a,a^2\}$ and labels $(1,1,0)$. The residual is
$g=-\mathbf1_{\{a^2\}}$. The amended series equals $1$ except at
$a^2$, where it equals $1-t$. Its $4\times4$ Hankel block has determinant
$t^3$, and its direct-sum construction has four states. Thus for every
$t\ne0$ its minimal dimension is exactly four, although the observed
$2\times2$ block only certifies dimension at least two. This exhibits the
gap between a rigorous lower certificate and a conservative but potentially
wasteful constructive upper bound.

\subsection{Evidence refinement produces quotients as well as inclusions}

There is also a categorical reason that capacity change cannot simply be
modeled as a canonical chain of embeddings. For a real behavior $f$, put
\[
 h_{S,E}:\mathbb R^S\longrightarrow\mathbb R^E,
 \qquad e_s\longmapsto(f(se))_{e\in E},\qquad
 Q_{S,E}=\operatorname{im}h_{S,E}.
\]
These are precisely finite observational image spaces of the kind used in
linear instances of CALF \cite{ce:calf2017}.

\begin{proposition}[Canonical comparison span and obstruction to a lift]
\label{ce:refinement-span}
For $S\subseteq S'$ and $E\subseteq E'$, set
$K=\operatorname{im}h_{S,E'}$. There is a canonical span
\[
 Q_{S,E}\ \xleftarrow{\ p\ }\ K\ \xrightarrow{\ i\ }\ Q_{S',E'},
\]
where $p$ is surjective coordinate restriction and $i$ is inclusion.
For a pure suffix refinement ($S'=S$), an old-to-new linear map
$j:Q_{S,E}\to Q_{S,E'}$ satisfying
$j h_{S,E}=h_{S,E'}$ exists if and only if
\[
 \ker h_{S,E}\subseteq\ker h_{S,E'}.
\]
Because the opposite kernel inclusion always holds, such a compatible lift
cannot exist when the suffix refinement strictly increases rank.
\end{proposition}
\begin{proof}
Restriction from $E'$ to $E$ takes every new observed row of an old prefix
to its old observed row. It therefore maps $K$ onto $Q_{S,E}$.
Adding prefixes includes its image in $Q_{S',E'}$. A linear map through
the quotient by $\ker h_{S,E}$ exists precisely when $h_{S,E'}$
annihilates that kernel. Finally, coordinate restriction implies
$\ker h_{S,E'}\subseteq\ker h_{S,E}$.
\end{proof}

For example, take the unary behavior with values
$f(\varepsilon)=f(a)=1$ and $f(a^k)=0$ for $k\geq2$.
With $S=\{\varepsilon,a\}$ and $E=\{\varepsilon\}$, both observed
rows equal $1$. Adding the suffix $a$ produces rows $(1,1)$ and $(1,0)$.
No function, let alone linear map, can send their single common old state
to both new states while respecting the represented prefixes.
Finite-dimensional surjections admit linear sections, but choosing a
section does not restore the impossible prefix compatibility.

There is a natural backward operation: the dual map
$p^*:Q_{S,E}^*\to K^*$ pulls old covectors back. Its image annihilates
$\ker p$, so it cannot supply the new error information in the directions
that newly added tests distinguish. Those directions require new evidence
and an initialization or refinement rule, not merely reverse transport of
old gradients. The span above is a comparison construction; there is no
claim that ordinary pullback composition of such spans automatically
recovers the span of a composite evidence refinement. Retaining generators
and observational presentations may be necessary for a compositional
framework.


\subsection{A model-language obstruction: unbounded rank with one real state}
\label{ce:nonlinear-scope}

\begin{proposition}[Linear Hankel dimension is not intrinsic state dimension]
\label{ce:hilbert}
The unary behavior $f(a^j)=1/(j+1)$ has no finite-dimensional real linear
WFA realization, yet has a one-dimensional smooth nonlinear realization on
the positive state domain:
\[
 x_0=1,\qquad x_{j+1}=\frac{x_j}{1+x_j},\qquad f(a^j)=x_j.
\]
Equivalently, in the coordinate $z=1/x$, the state update is
$z_{j+1}=z_j+1$ and the output is $1/z_j$.
\end{proposition}
\begin{proof}
Induction gives $x_j=1/(j+1)$. Its $m\times m$ Hankel block, with
indices $0,\ldots,m-1$, has entries $H_{i,j}=1/(i+j+1)$. For every
nonzero $v\in\mathbb R^m$,
\[
 v^THv=\int_0^1\left(\sum_{i=0}^{m-1}v_i t^i\right)^2\,dt>0,
\]
since a nonzero real polynomial is nonzero on an open subinterval.
Thus every such block has rank $m$, excluding all finite linear WFA
realizations. The coordinate change is smooth on the positive domain and
transforms the displayed update as claimed.
\end{proof}

The sequence itself is not a new separating example; Gastin and Monmege
use $1/(|u|+1)$ in Remark~2 of their study of weighted automata
\cite{ce:gastin2014}, in a different comparison of model classes. The
argument above directly proves the required real-linear rank statement.
Nonlinear weighted automata are also an existing formalism
\cite{ce:nonlinear2018}.

This countermodel materially changes the interpretation of a rank-growth
policy. Ever more WFA states can reflect a mismatch in the allowed
transition or output language rather than an increasingly complicated
latent mechanism. A behavioral rank certificate is relative to its chosen
representation class. It is not a domain-general explanation of complexity,
and changing coordinates cannot make its conclusion stronger than that.
No finite exact word dataset can refute the entire class of finite WFAs:
the finite-support trie construction always supplies a fitting finite WFA.
Longer horizons can refute a fixed dimension budget, but cannot turn finite
non-refutation into a proof of the full target language.

Nor can word observations distinguish implementations with identical word
behavior. Additional interventions, alternative initial conditions, or new
observation channels can only be invoked when the experimental setting
actually permits them. The objective of predicting word outputs, the
objective of learning a minimal linear realization, and the objective of
identifying a physical or cognitive mechanism are different stakeholder
objectives, even when they use some of the same data.

\subsection{What has been obtained, and what would constitute a next theorem}

The results give a rigorous, bounded integration rather than a universal
identification of learning and reconstruction:
\begin{enumerate}
\item \emph{Evidence:} a finite Hankel singular value above a justified
      noise floor certifies insufficient behavioral state dimension.
\item \emph{Within a matrix relaxation:} after fixed-rank stationarity,
      the certificate forces a normal innovation. Its leading singular
      vectors produce a rank-one, worst-case optimal noise-robust birth
      with an explicit true-table descent guarantee.
\item \emph{Within actual automata:} a finite residual has an explicit
      realization. A gated direct sum preserves all old behavior at
      initialization, preserves old gradients, and creates a live new
      gradient. It is a constructive upper bound, not a minimal solution.
\item \emph{Across parameter spaces:} reverse-mode compatibility is a
      pullback law; exact optimizer commutation would prohibit the normal
      directions needed for escape.
\item \emph{Across evidence spaces:} adding generators gives inclusions,
      while adding tests refines quotients. There is no general canonical
      old-state embedding consistent with newly distinguished histories.
\end{enumerate}

The substantive next problem is to close the demonstrated gap: construct
\emph{small, Hankel-consistent, transition-consistent residual extensions}
whose data-selected gates enjoy a robust descent bound, and identify the
conditions under which their dimension is controlled by the observed
innovation rather than by a memorizing trie. The unconstrained leading
singular vector does not solve this, and the explicit counterexample shows
why. This target is more concrete than asserting a universal common
reconstruction theorem for automata, formal concepts, and topology: it asks
for an actual controlled operation and a theorem that fails without its
structural hypotheses.

A principled first enlargement of scope is from real WFAs to linear
second-order recurrent networks, where a precise spectral equivalence is
already available \cite{ce:rabusseau2019}. The WFA-specific Hankel rank
certificate does not automatically measure the intrinsic dimension of a
general nonlinear recurrent state space. Likewise, rank-adaptive matrix or
tensor training already exists \cite{ce:tan2014,ce:lottery2022}; claiming
that adaptive capacity alone is new would be incorrect.

\paragraph{Primary-source positioning.}
CALF supplies the categorical image, closedness, and consistency viewpoint
\cite{ce:calf2017}. Parametric lenses supply compositional reverse-mode and
optimizer semantics \cite{ce:lenses2024}. Net2Net supplies prior examples
of function-preserving architecture changes \cite{ce:net2net2016}.
Riemannian pursuit already alternates fixed-rank optimization and rank
increase \cite{ce:tan2014}; GradMax explicitly preserves old network
outputs while choosing new weights through gradient maximization and an
SVD \cite{ce:gradmax2022}; dynamical low-rank methods already adapt ranks
during neural training \cite{ce:lottery2022}. Constrained Hankel completion
and Hankel concentration address two essential statistical/structural
issues \cite{ce:balle2012,ce:denis2016}. The robust minimax calculation,
the explicit dead/live gate comparison, and the refinement-span analysis
above are proved here so their exact scope can be checked; they are not
presented as a literature-verified novelty claim.

\begin{thebibliography}{99}
\bibitem{ce:calf2017}
G. van Heerdt, M. Sammartino, and A. Silva.
\emph{CALF: Categorical Automata Learning Framework}. CSL 2017,
LIPIcs 82, 29:1--29:24.
\url{https://doi.org/10.4230/LIPIcs.CSL.2017.29}.

\bibitem{ce:lenses2024}
G. S. H. Cruttwell, B. Gavranovic, N. Ghani, P. Wilson, and F. Zanasi.
\emph{Deep Learning with Parametric Lenses}. 2024.
\url{https://arxiv.org/abs/2404.00408}.

\bibitem{ce:tan2014}
M. Tan, I. W. Tsang, L. Wang, B. Vandereycken, and S. J. Pan.
\emph{Riemannian Pursuit for Big Matrix Recovery}.
ICML 2014, PMLR 32(2), 1539--1547.
\url{https://proceedings.mlr.press/v32/tan14.html}.

\bibitem{ce:net2net2016}
T. Chen, I. Goodfellow, and J. Shlens.
\emph{Net2Net: Accelerating Learning via Knowledge Transfer}.
2015 preprint, revised 2016.
\url{https://arxiv.org/abs/1511.05641}.

\bibitem{ce:balle2012}
B. Balle and M. Mohri.
\emph{Spectral Learning of General Weighted Automata via Constrained
Matrix Completion}. NeurIPS 25, 2012.
\url{https://papers.neurips.cc/paper_files/paper/2012/hash/700fdb2ba62d4554dc268c65add4b16e-Abstract.html}.

\bibitem{ce:denis2016}
F. Denis, M. Gybels, and A. Habrard.
\emph{Dimension-free Concentration Bounds on Hankel Matrices for Spectral
Learning}. JMLR 17(31), 1--32, 2016.
\url{https://jmlr.org/papers/v17/14-501.html}.

\bibitem{ce:rabusseau2019}
G. Rabusseau, T. Li, and D. Precup.
\emph{Connecting Weighted Automata and Recurrent Neural Networks through
Spectral Learning}. AISTATS 2019, PMLR 89, 1630--1639.
\url{https://proceedings.mlr.press/v89/rabusseau19a.html}.

\bibitem{ce:lottery2022}
S. Schotth\"ofer, E. Zangrando, J. Kusch, G. Ceruti, and F. Tudisco.
\emph{Low-rank lottery tickets: finding efficient low-rank neural networks
via matrix differential equations}. 2022.
\url{https://arxiv.org/abs/2205.13571}.

\bibitem{ce:gradmax2022}
U. Evci, B. van Merri\"enboer, T. Unterthiner, M. Vladymyrov, and F. Pedregosa.
\emph{GradMax: Growing Neural Networks using Gradient Information}.
ICLR 2022. \url{https://arxiv.org/abs/2201.05125}.

\bibitem{ce:gastin2014}
P. Gastin and B. Monmege.
\emph{Adding Pebbles to Weighted Automata: Easy Specification and Efficient
Evaluation?} Theoretical Computer Science, 2014. In particular Remark~2.
\url{https://lsv.ens-paris-saclay.fr/Publis/PAPERS/PDF/GM-tcs14.pdf}.

\bibitem{ce:nonlinear2018}
T. Li, G. Rabusseau, and D. Precup.
\emph{Nonlinear Weighted Finite Automata}.
AISTATS 2018, PMLR 84, 679--688.
\url{https://proceedings.mlr.press/v84/li18a.html}.

\end{thebibliography}

\end{document}
